% ============================================================
% Chapter 7 — Level 02: Authentication Security
% ============================================================
\chapter{Level 02 — Authentication Security}
\label{ch:level02}

\noteblock{This module is a \textbf{placeholder}. The lab environment for Level 02 (\texttt{auth.lab.uzyntra.com}) will be deployed separately. This chapter documents the learning objectives, vulnerability theory, and testing methodology that will be used when the lab is live.}

% ─────────────────────────────────────────────────────────
\section{Learning Objectives}

By completing this module you will be able to:
\begin{itemize}
  \item Identify and exploit weak login security controls
  \item Analyze and attack session management implementations
  \item Test OTP (One-Time Password) workflows for vulnerabilities
  \item Evaluate MFA (Multi-Factor Authentication) implementation quality
  \item Use Burp Suite Intruder and Repeater for authentication attacks
\end{itemize}

% ─────────────────────────────────────────────────────────
\section{Vulnerability Explanations}

\subsection{Login Security}

Authentication vulnerabilities allow attackers to bypass or subvert the login process.

\begin{table}[H]
\centering
\caption{Common login vulnerabilities}
\begin{tabularx}{\textwidth}{>{\bfseries}lXl}
\toprule
Vulnerability & Description & CWE \\
\midrule
Username enumeration & Different response for valid vs invalid username & CWE-204 \\
No rate limiting & Unlimited login attempts permitted & CWE-307 \\
Weak password policy & Short or simple passwords accepted & CWE-521 \\
Default credentials & admin/admin, admin/password accepted & CWE-1392 \\
Verbose errors & ``Password incorrect'' vs ``User not found'' & CWE-209 \\
\bottomrule
\end{tabularx}
\end{table}

\subsection{Session Management}

After authentication, the application issues a session token. Weak session management allows token theft or prediction.

\begin{table}[H]
\centering
\caption{Session management vulnerabilities}
\begin{tabularx}{\textwidth}{>{\bfseries}lXl}
\toprule
Vulnerability & Description & CWE \\
\midrule
Session fixation & Attacker sets a known token before login & CWE-384 \\
Predictable tokens & Sequential or low-entropy session IDs & CWE-330 \\
No token rotation & Token unchanged after privilege change & CWE-613 \\
Missing cookie flags & No HttpOnly, Secure, or SameSite & CWE-614 \\
Long session lifetime & Token valid for weeks after logout & CWE-613 \\
\bottomrule
\end{tabularx}
\end{table}

\subsection{OTP Workflow Testing}

OTP vulnerabilities allow attackers to bypass two-factor authentication.

\begin{itemize}
  \item \textbf{OTP not expiring} — codes valid indefinitely
  \item \textbf{OTP reuse} — same code accepted multiple times
  \item \textbf{OTP brute force} — 6-digit code: only 1,000,000 possibilities
  \item \textbf{OTP bypass} — application accepts login without OTP step
  \item \textbf{Predictable OTP} — TOTP seed reused or derived from username
\end{itemize}

\subsection{MFA Implementation Review}

\begin{itemize}
  \item Is MFA enforced for all admin accounts?
  \item Can MFA be skipped by directly navigating to post-login URLs?
  \item Is the MFA backup code single-use and properly invalidated?
  \item Does the application correctly validate the TOTP window?
\end{itemize}

% ─────────────────────────────────────────────────────────
\section{Required Lab Setup}

\begin{table}[H]
\centering
\caption{Level 02 Lab Configuration (Pending Deployment)}
\begin{tabular}{ll}
\toprule
Field & Value \\
\midrule
Target & \texttt{https://auth.lab.uzyntra.com} (pending) \\
Test accounts & \texttt{user@lab.uzyntra.com / training123} \\
Admin account & \texttt{admin@lab.uzyntra.com / admin123} \\
Burp proxy & \texttt{127.0.0.1:8080} \\
\bottomrule
\end{tabular}
\end{table}

% ─────────────────────────────────────────────────────────
\section{Burp Suite Testing Process}

\subsection{Testing for Username Enumeration}

\begin{enumerate}
  \item Capture the login request in Burp Proxy
  \item Send to Repeater
  \item Try a known valid username with wrong password $\to$ note response length/message
  \item Try an invalid username with wrong password $\to$ compare response
  \item If responses differ $\to$ username enumeration confirmed
\end{enumerate}

\subsection{Testing for Brute Force (Intruder)}

\begin{enumerate}
  \item Capture login POST request $\to$ Send to Intruder
  \item Set position on password field: \texttt{password=§test§}
  \item Load payload: \texttt{/usr/share/seclists/Passwords/Common-Credentials/10k-most-common.txt}
  \item Start attack $\to$ sort by response length or status code
  \item Different length/code = successful login = weak password confirmed
\end{enumerate}

\warnblock{In Community Edition, Intruder is throttled. For brute force testing in real engagements, use \texttt{hydra} or \texttt{ffuf} for speed. Intruder is fine for learning the concept.}

\subsection{Testing Session Token Entropy}

\begin{enumerate}
  \item Log in 5 times with the same account
  \item Capture each Set-Cookie header from the login response
  \item Send all 5 cookies to Sequencer
  \item Run the analysis — Burp evaluates token randomness
  \item Low entropy = predictable tokens = session hijacking risk
\end{enumerate}

% ─────────────────────────────────────────────────────────
\section{Security Testing Exercise}

\exerciseblock{
\textbf{Exercise 7.1 — Authentication Reconnaissance (Theory)}\\[4pt]
While the lab environment is being set up, practice the analysis:
\begin{enumerate}
  \item Describe what HTTP response differences would confirm username enumeration
  \item Write a Burp Intruder payload position for: \texttt{POST /login} with body \texttt{email=admin@test.com\&password=FUZZ}
  \item List 5 cookie flag configurations that indicate a security weakness
  \item Describe how you would test if an OTP code can be reused
\end{enumerate}
}

% ─────────────────────────────────────────────────────────
\section{Professional Finding Report Template}

\findingblock{
\begin{tabular}{ll}
\textbf{Finding Title:} & Username Enumeration via Login Response \\
\textbf{Severity:} & \textcolor{uzyellow}{\textbf{Medium}} \\
\textbf{CWE:} & CWE-204 \\
\textbf{OWASP:} & A07:2021 — Identification and Authentication Failures \\
\textbf{Affected Endpoint:} & \texttt{POST /api/auth/login} \\
\end{tabular}\\[8pt]
\textbf{Description:}\\
The login endpoint returns different HTTP responses for valid vs invalid usernames, allowing an attacker to enumerate valid user accounts.\\[6pt]
\textbf{Evidence:}\\
Valid user, wrong password: HTTP 200, body: \texttt{"error":"Invalid password"}\\
Invalid user: HTTP 200, body: \texttt{"error":"User not found"}\\[6pt]
\textbf{Impact:}\\
Attacker can enumerate all valid email addresses and target them for brute force or phishing.\\[6pt]
\textbf{Recommendation:}\\
Return identical responses for all failed login attempts: \texttt{"Invalid credentials"}.
}

% ─────────────────────────────────────────────────────────
\section{Module Completion Checklist}

\begin{tcolorbox}[findingstyle, title={Level 02 Burp Checklist}]
\begin{itemize}
  \item[\checkbox] Login request captured in Burp HTTP History
  \item[\checkbox] Username enumeration tested (valid vs invalid user responses compared)
  \item[\checkbox] Rate limiting tested (50+ login attempts — any lockout?)
  \item[\checkbox] Password policy tested (accepted: \texttt{a}, \texttt{123}, common passwords?)
  \item[\checkbox] Session token captured from Set-Cookie header
  \item[\checkbox] Cookie flags checked: HttpOnly, Secure, SameSite
  \item[\checkbox] Session token analyzed in Sequencer for entropy
  \item[\checkbox] Token invalidated after logout (send old token post-logout)
  \item[\checkbox] OTP reuse tested (same code submitted twice)
  \item[\checkbox] OTP brute force rate limiting tested
  \item[\checkbox] MFA skip tested (navigate directly to /dashboard after step 1)
  \item[\checkbox] All findings documented with request/response evidence
\end{itemize}
\end{tcolorbox}
